\documentclass[spanish,letterpaper,oneside]{article}

\def\documenttitle {Del Diagnóstico Práctico a una Propuesta de Innovación Matemática}
\def\documentsubtitle {Fundamentos para la enseñanza y el aprendizaje I · Cohorte SD26}
\def\documentsubject {Actividad 6 · Bloque III}
\def\documentauthor {Mtro. Martín Jonathan de la Cruz Muñoz}
\def\coursename {Fundamentos para la enseñanza y el aprendizaje I}
\def\coursecode {MAEM-SD26}

\def\universityname {Instituto de Investigación, Innovación y Estudios de Posgrado para la Educación del Estado de Nuevo León}
\def\universityfaculty {Maestría en Aprendizaje y Enseñanza de las Matemáticas}
\def\universitydepartment {Fundamentos para la enseñanza y el aprendizaje I}
\def\universitydepartmentimage {img/iiiepe/marca-iiiepe.png}
\def\universitydepartmentimagecfg {height=1.75cm}
\def\universitylocation {Monterrey, Nuevo León}
\def\studentfullname {Mtro. Martín Jonathan de la Cruz Muñoz}
\def\studentid {00874}
\def\studentemail {martin.delacruz@campus.iiiepe.edu.mx}
\def\programperiod {2do periodo, septiembre--diciembre de 2026 (SD26)}
\def\programmodality {Modalidad presencial}
\def\programshort {MAEM}

\def\authortable {
  \begin{tabular}{ll}
    \textbf{Alumno:} & \studentfullname \\
    Matrícula: & \studentid \\
    Correo institucional: & \studentemail \\
    Unidad: & \coursename \\
    Programa: & \programshort \\
    Periodo: & \programperiod \\
    & \\
    \multicolumn{2}{l}{Fecha de realización: 19 de noviembre de 2026} \\
    \multicolumn{2}{l}{\universitylocation}
  \end{tabular}
}

\input{template}

\usepackage{helvet}
\renewcommand{\familydefault}{\sfdefault}
\usepackage{array}
\usepackage{longtable}
\usepackage{pdflscape}
\usepackage{tikz}
\usepackage{pgfplots}
\pgfplotsset{compat=1.18}
\usepackage[most]{tcolorbox}
\usetikzlibrary{arrows.meta,positioning,calc,fit,backgrounds,shapes.geometric}

\definecolor{iiiepeNavy}{HTML}{163B5C}
\definecolor{iiiepeBlue}{HTML}{1F6F9F}
\definecolor{iiiepeAqua}{HTML}{20B7A8}
\definecolor{iiiepeOrange}{HTML}{F28C28}
\definecolor{iiiepePaper}{HTML}{F6F8F7}
\definecolor{iiiepeInk}{HTML}{1E2A32}
\definecolor{iiiepePale}{HTML}{EEF5F8}
\definecolor{iiiepeMint}{HTML}{E7F6F4}
\definecolor{iiiepeWarm}{HTML}{FFF3E6}
\def\portraittitlecolor {iiiepeNavy}
\setcitestyle{authoryear,open={(},close={)}}
\makeatletter
\renewcommand{\@biblabel}[1]{}
\makeatother
\setlength{\bibhang}{1.27cm}
\setlength{\bibsep}{6pt}
\AtBeginDocument{\renewcommand{\refname}{Referencias}}

\newcommand{\insertcoverwatermark}{%
  \AddToShipoutPictureBG*{%
    \begin{tikzpicture}[remember picture,overlay]
      \node[opacity=0.04,inner sep=0pt] at (current page.center) {%
        \includegraphics[width=0.65\paperwidth]{img/iiiepe/marca-iiiepe.png}%
      };
    \end{tikzpicture}%
  }%
}

\setlength{\parskip}{5pt}\setlength{\parindent}{0pt}
\begin{document}\insertcoverwatermark\templatePortrait\templatePagecfg\templateFinalcfg\fancyfoot{}\fancyfoot[L]{\scriptsize IIIEPE · MAEM · SD26}\fancyfoot[C]{\scriptsize\thepage}\fancyfoot[R]{\scriptsize Actividad 6 · Tercer Momento}\setstretch{1.3}\raggedbottom
\section{Introducción y postura central}
La transición hacia una enseñanza matemática significativa demanda articular el marco conceptual del aprendizaje profundo con un diagnóstico honesto de la propia práctica y una propuesta concreta de innovación didáctica. Este estudio integra de manera orgánica tres momentos progresivos de construcción pedagógica que conducen desde la reflexión teórica hasta el diseño de una intervención innovadora. Mi postura es que las dificultades matemáticas más persistentes en educación media superior técnica e ingeniería no surgen de una incapacidad intrínseca del estudiante, sino del predominio de prácticas didácticas fragmentadas que separan el cálculo simbólico de la modelación, la colaboración entre pares y la reflexión metacognitiva \citep{quiroga2022aprendizaje,sawyer2008optimizing}.

Considero necesario fundamentar cada decisión de cambio en evidencias reales de enseñanza (producciones de cuadernos, guías resueltas y registros evaluativos) sin responsabilizar unilateralmente al alumnado. Por ello, tras presentar la síntesis analítica de \citet{quiroga2022aprendizaje} (Primer momento) y el cuadro diagnóstico de seis dimensiones (Segundo momento), se expone el diseño completo de una propuesta de innovación \textbf{por realizar} (Tercer momento) orientada a transformar el abordaje de funciones lineales y variación afín en bachillerato técnico e ingeniería \citep{rivas2017innovaciones,conalep2025mecatronica}.

\section{Primer momento: análisis y reflexión de la lectura base}
A partir de la lectura analítica de \citet{quiroga2022aprendizaje} y la conceptualización de las nuevas pedagogías, el examen de cinco ejes articuladores permite esclarecer el sentido pedagógico de innovar:
\begin{enumerate}
  \item \textbf{Diferencias entre aprendizaje superficial y profundo:} Recuperando a \citet{marton1976qualitative}, el enfoque superficial busca reproducir fórmulas o definiciones para resolver la tarea inmediata, generando saberes inertes; en cambio, el aprendizaje profundo conecta conceptos previos con nuevos significados, argumenta con evidencia y transfiere estrategias hacia problemas reales \citep{sawyer2008optimizing}.
  \item \textbf{Dimensiones cognitiva, interpersonal e intrapersonal:} La comprensión matemática duradera requiere activar simultáneamente procesos cognitivos de análisis y modelación, interacciones interpersonales de debate y comunicación clara en equipo, así como disposiciones intrapersonales de autorregulación, perseverancia y mentalidad de crecimiento \citep{quiroga2022aprendizaje,schunk2012teorias}.
  \item \textbf{Transformación en los roles docente y estudiantil:} Retomando a \citet{fullan2014rich}, las nuevas pedagogías desplazan el control unidireccional docente hacia relaciones de \textit{asociación} donde profesores y estudiantes colaboran como co-aprendices, pasando gradualmente hacia la autonomía estudiantil guiada.
  \item \textbf{Características de las nuevas pedagogías:} Priorizan tareas auténticas vinculadas al entorno, creación y experimentación de conocimiento en lugar de recepción pasiva, uso intencionado de tecnología para ampliar la indagación y prácticas de liderazgo pedagógico orientadas por metas compartidas \citep{rivas2017innovaciones}.
  \item \textbf{Creencias pedagógicas y cultura escolar:} Cambiar actividades en superficie no transforma el aula si permanecen intactas las creencias sobre qué define una buena clase de matemáticas o un «buen alumno». La cultura escolar solo evoluciona cuando se interviene conjuntamente sobre recursos, didáctica y significados pedagógicos compartidos \citep{quiroga2022aprendizaje}.
\end{enumerate}

\section{Segundo momento: análisis diagnóstico de la propia práctica de enseñanza}
El análisis propio se sitúa en una problemática recurrente observada en grupos de cuarto semestre de bachillerato técnico (mantenimiento industrial y mecatrónica) y primer ingreso de ingeniería al abordar \textbf{modelación con funciones lineales y comparación entre variación afín y proporcionalidad}: al revisar cuadernos de clase, guías de ejercicios y listas de cotejo de evaluación parcial, se advierte que el estudiantado ejecuta despejes correctos cuando recibe ecuaciones preestablecidas, pero fracasa al formular el modelo desde una gráfica experimental o al interpretar qué significan pendiente, ordenada al origen, unidades físicas y dominio de validez \citep{star2015algebra}. Evitando atribuir este hallazgo únicamente a carencias del estudiantado, el cuadro siguiente examina las seis dimensiones críticas del diagnóstico docente:

{\small\setstretch{1.15}\renewcommand{\arraystretch}{1.35}
\begin{longtable}{|>{\raggedright\arraybackslash}p{0.20\linewidth}|>{\raggedright\arraybackslash}p{0.25\linewidth}|>{\raggedright\arraybackslash}p{0.49\linewidth}|}\hline
\rowcolor{iiiepePale} \textbf{Aspecto diagnóstico} & \textbf{Pregunta orientadora} & \textbf{Análisis situado con evidencias concretas} \\\hline
\endfirsthead
\hline
\rowcolor{iiiepePale} \textbf{Aspecto diagnóstico} & \textbf{Pregunta orientadora} & \textbf{Análisis situado con evidencias concretas} \\\hline
\endhead
1. Problemática identificada & ¿Qué dificultades presentan los alumnos al aprender matemáticas? & En 18 de aproximadamente 25 cuadernos revisados de secuencias previas sobre rectas, los estudiantes obtienen $m$ y $b$ por sustitución en fórmula, pero no explican su significado físico (e.g., tasa de variación frente a valor inicial ni intervalo real de operación) ni identifican si $y=kx$ difiere de $y=mx+b$. \\\hline
2. Práctica de enseñanza actual & ¿Cómo abordo la enseñanza de ese contenido matemático? & Inicio frecuentemente definiendo $y=mx+b$ en el pizarrón, resuelvo dos casos tipo con pasos cerrados y asigno guías individuales de práctica operativa antes de abordar aplicaciones técnicas. \\\hline
3. Concepciones pedagógicas & ¿Qué ideas sobre enseñar y aprender matemáticas orientan mi práctica? & Subyace la creencia heredada de la formación ingenieril de que dominar primero el mecanismo algebraico habilita después la aplicación contextual, subestimando que el sentido físico construye el concepto algebraico. \\\hline
4. Características del aprendizaje & ¿Qué evidencias muestran un aprendizaje superficial o profundo? & Superficial: repetición correcta del algoritmo en hojas con formato fijo pero bloqueo cuando cambia el nombre de la variable o se solicita interpretación escrita. Profundo incipiente: preguntas espontáneas de taller cuando conectan pendiente con sensibilidad de un sensor real. \\\hline
5. Causas posibles & ¿Qué decisiones didácticas podrían estar contribuyendo a la problemática? & Privilegiar ejercicios sin contexto en los primeros bloques, resolver dudas demasiado rápido en el pizarrón por presión del horario y ponderar casi únicamente resultados numéricos en las evaluaciones escritas \citep{eef2017mathematics}. \\\hline
6. Necesidad de transformación & ¿Qué tendría que modificar en mi práctica para favorecer el aprendizaje profundo? & Invertir la secuencia: partir de un reto técnico auténtico con datos, organizar colaboración deliberada con roles de argumentación, implementar andamiaje gradual y evaluar mediante matrices formativas que valoren explicación y transferencia \citep{eef2021feedback}. \\\hline
\end{longtable}}

\section{Tercer momento: propuesta de innovación didáctica (por realizar)}
Con base en el diagnóstico previo y el marco de \citet{quiroga2022aprendizaje} y \citet{rivas2017innovaciones}, diseño la siguiente propuesta de innovación pedagógica —formulada como \textbf{intervención por realizar}— en el marco del desarrollo profesional en la MAEM:

\begin{tcolorbox}[colback=iiiepePale,colframe=iiiepeNavy,title=\textbf{Ficha técnica de la innovación · «Calibrar y comparar para decidir» (por realizar)},breakable]
\begin{enumerate}
  \item \textbf{Contenido matemático (tema y nivel):} Modelación lineal ($y = mx + b$), distinción frente a proporcionalidad directa ($y = kx$), interpretación de tasa de cambio, ordenada al origen, punto de equilibrio e intervalo de validez, en cuarto semestre de bachillerato técnico en mecatrónica y mantenimiento (adaptable a cálculo inicial de ingeniería) \citep{sep2023pensamiento,conalep2025mecatronica}.
  \item \textbf{Problemática atendida:} Manipulación algebraica mecánica desconectada del significado físico, dimensional y argumentativo de los parámetros lineales.
  \item \textbf{Propósito de aprendizaje:} Lograr que el estudiantado construya, contraste y justifique modelos lineales desde tablas de datos experimentales o simulados, articulando representaciones numérica, gráfica y algebraica para sustentar decisiones técnicas con evidencia.
  \item \textbf{Situación o desafío matemático:} Dos dispositivos de monitoreo térmico de banco ($S_1$ calibrado sin sesgo con razón $1.00$ y $S_2$ con lectura simulada $y = 1.05x + 0.2$ respecto a patrón $x$) y dos tarifas de mantenimiento preventivo ($C_A(h) = 250 + 80h$ y $C_B(h) = 100 + 110h$) presentan comportamientos distintos. Con datos didácticos simulados suministrados en hoja de taller, cada equipo debe responder: \textit{¿en qué rango de operación resulta confiable $S_2$, qué error introduce asumir proporcionalidad directa y a partir de qué duración conviene cambiar de esquema de servicio?} Deben predecir antes de calcular, validar unidades y defender su dictamen por escrito.
  \item \textbf{Estrategias de enseñanza:} Secuencia de tres sesiones basada en asociación docente--estudiante \citep{fullan2014rich}: (a) apertura con predicción individual sin fórmula para activar saberes previos; (b) exploración en equipos de tres con roles rotativos (\textit{modelador gráfico}, \textit{verificador dimensional} y \textit{portavoz argumentativo}); (c) retirada gradual de guías de preguntas hacia autonomía progresiva; y (d) debate plenario estructurado sobre errores productivos y límites del modelo.
  \item \textbf{Integración de las dimensiones del aprendizaje profundo:}
    \begin{itemize}
      \item \textit{Cognitiva:} Contraste entre tabla de diferencias finitas, pendiente gráfica y ecuación simbólica; verificación de límites físicos.
      \item \textit{Interpersonal:} Negociación de hipótesis entre pares, coevaluación respetuosa entre equipos y construcción de consenso argumentado.
      \item \textit{Intrapersonal:} Diario breve de metacognición tras cada sesión («qué supuse al inicio, qué evidencia me hizo corregir el modelo y qué aprendí del error») para fomentar mentalidad de crecimiento \citep{schunk2012teorias}.
    \end{itemize}
  \item \textbf{Evaluación formativa (evidencias e indicadores):} Recolección continua de cuatro evidencias: predicción diagnóstica inicial, hoja de análisis tabular--gráfico por equipo, dictamen escrito justificado y reflexión metacognitiva individual revisada tras retroalimentación descriptiva \citep{bostrom2023formative}.
\end{enumerate}
\end{tcolorbox}

La Figura~\ref{fig:act6-secuencia} ilustra el flujo metodológico propuesto y el cuadro siguiente define los indicadores formativos de avance:

\begin{figure}[htbp]
\centering
\begin{tikzpicture}[font=\scriptsize, >=Latex, scale=0.92, transform shape]
  \node[draw=iiiepeNavy,fill=iiiepePale,thick,rounded corners=4pt,text width=2.8cm,minimum height=2.0cm,align=center] (f1) at (0,0) {\textbf{Sesión 1}\\Predicción inicial y análisis tabular\\(\textit{Saberes previos})};
  \node[draw=iiiepeBlue,fill=iiiepePale,thick,rounded corners=4pt,text width=2.8cm,minimum height=2.0cm,align=center] (f2) at (3.65,0) {\textbf{Sesión 2}\\Modelación gráfica y algebraica en tríos\\(\textit{Cognitivo + Social})};
  \node[draw=iiiepeAqua!80!black,fill=iiiepeMint,thick,rounded corners=4pt,text width=2.8cm,minimum height=2.0cm,align=center] (f3) at (7.3,0) {\textbf{Sesión 3}\\Debate de validez y retroalimentación\\(\textit{Evaluación formativa})};
  \node[draw=iiiepeOrange!90!black,fill=iiiepeWarm,thick,rounded corners=4pt,text width=2.8cm,minimum height=2.0cm,align=center] (f4) at (10.95,0) {\textbf{Cierre}\\Segunda versión y diario metacognitivo\\(\textit{Intrapersonal})};

  \draw[->,very thick,iiiepeNavy] (f1) -- (f2);
  \draw[->,very thick,iiiepeBlue] (f2) -- (f3);
  \draw[->,very thick,iiiepeAqua!80!black] (f3) -- (f4);
\end{tikzpicture}
\caption{Secuencia didáctica de innovación en tres sesiones articulando las dimensiones del aprendizaje profundo.}
\label{fig:act6-secuencia}
\end{figure}

{\small\setstretch{1.15}\renewcommand{\arraystretch}{1.3}
\begin{longtable}{|>{\raggedright\arraybackslash}p{0.22\linewidth}|>{\raggedright\arraybackslash}p{0.34\linewidth}|>{\raggedright\arraybackslash}p{0.38\linewidth}|}\hline
\rowcolor{iiiepePale} \textbf{Dimensión evaluada} & \textbf{Evidencia observable} & \textbf{Indicador formativo de avance} \\\hline
\endfirsthead
\hline
\rowcolor{iiiepePale} \textbf{Dimensión evaluada} & \textbf{Evidencia observable} & \textbf{Indicador formativo de avance} \\\hline
\endhead
Comprensión conceptual y modelación & Hojas de registro tabular--gráfico y ecuación justificada. & Distingue relación proporcional vs afín e interpreta pendiente y ordenada al origen con unidades físicas correctas. \\\hline
Argumentación y transferencia & Dictamen técnico escrito comparando intervalos y umbrales. & Justifica decisiones bajo condiciones cambiantes señalando al menos un límite físico o supuesto del modelo. \\\hline
Colaboración e intrapersonalidad & Observación de roles en equipo y diario de corrección tras devolución. & Recupera argumentos de compañeros, reconoce obstáculos iniciales y mejora explícitamente su segunda entrega \citep{eef2021feedback}. \\\hline
\end{longtable}}

\section{Transferencia profesional e impacto esperado en la práctica docente}
Esta propuesta impacta directamente la coherencia entre formación técnica e ingeniería y enseñanza comprensiva de las matemáticas. Además, permite articular los contenidos de bachillerato tecnológico con problemas genuinos de medición; sin embargo, su eficacia dependerá de sostener la retroalimentación formativa a lo largo de las tres sesiones. Dado que se plantea como proyecto por implementar y sistematizar mediante ciclos de investigación-acción \citep{restrepo2004investigacion}, sus evidencias futuras permitirán verificar empíricamente en qué medida el desplazamiento desde la ejercitación mecánica hacia la indagación colaborativa fortalece la autonomía y la argumentación crítica en el plantel. Es decir, por ello, no obstante los límites propios de un horario escolar acotado, el diseño prioriza evidencias comprensivas sobre la cobertura rutinaria.

\clearpage
\section{Conclusiones y síntesis argumentativa}
En conclusión, el desarrollo articulado de la fundamentación teórica, el diagnóstico de la práctica y el diseño de innovación confirma que transformar la enseñanza de las matemáticas exige mirar primero hacia nuestras propias decisiones y creencias antes de diagnosticar carencias en los estudiantes. El análisis crítico de \citet{quiroga2022aprendizaje} aporta la brújula conceptual necesaria para articular pensamiento profundo, asociación interpersonal y cultura escolar colaborativa.

Por tanto, el diseño de innovación presentado traduce dichos principios en una secuencia operativa realista, éticamente responsable y verificable mediante evaluación formativa continua, manteniendo coherencia plena con el tránsito docente humanista que consolido en la maestría.

\clearpage
\begin{thebibliography}{99}
\bibitem[Boström y Palm(2023)]{bostrom2023formative} Boström, E., y Palm, T. (2023). Formative assessment practice in mathematics: Teachers' perceptions of prerequisite conditions. \textit{Frontiers in Education}, 8, Artículo 1101192.
\bibitem[Collin y Quigley(2021)]{eef2021feedback} Collin, J., y Quigley, A. (2021). \textit{Teacher feedback to improve pupil learning}. Education Endowment Foundation.
\bibitem[CONALEP(2025)]{conalep2025mecatronica} Colegio Nacional de Educación Profesional Técnica. (2025). \textit{Perfil de egreso: Profesional Técnico-Bachiller en Mecatrónica. Modelo Académico 2025}. Secretaría Académica.
\bibitem[Fullan y Langworthy(2014)]{fullan2014rich} Fullan, M., y Langworthy, M. (2014). \textit{A rich seam: How new pedagogies find deep learning}. Pearson.
\bibitem[Henderson et~al.(2017)]{eef2017mathematics} Henderson, P., Hodgen, J., Foster, C., y Kuchemann, D. (2017). \textit{Improving mathematics in Key Stages 2 and 3}. Education Endowment Foundation.
\bibitem[Marton y Säljö(1976)]{marton1976qualitative} Marton, F., y Säljö, R. (1976). On qualitative differences in learning: I. Outcome and process. \textit{British Journal of Educational Psychology}, 46(1), 4--11.
\bibitem[Quiroga Aguilar y Lara Órdenes(2022)]{quiroga2022aprendizaje} Quiroga Aguilar, L., y Lara Órdenes, E. (2022). El aprendizaje profundo como herramienta para cambio en la visión de aprendizaje de una cultura escolar. \textit{Revista Educación las Américas}, 12(1). \url{https://doi.org/10.35811/rea.v12i1.201}
\bibitem[Restrepo Gómez(2004)]{restrepo2004investigacion} Restrepo Gómez, B. (2004). La investigación-acción educativa y la construcción de saber pedagógico. \textit{Educación y Educadores}, (7), 45--55.
\bibitem[Rivas et~al.(2017)]{rivas2017innovaciones} Rivas, A., André, F., y Delgado, L. E. (Eds.). (2017). \textit{50 innovaciones educativas para escuelas}. CIPPEC / Santillana.
\bibitem[Sawyer(2008)]{sawyer2008optimizing} Sawyer, R. K. (2008). Optimizing learning: Implications of learning sciences research. OECD Publishing.
\bibitem[Schunk(2012)]{schunk2012teorias} Schunk, D. H. (2012). \textit{Teorías del aprendizaje: Una perspectiva educativa} (6.ª ed.). Pearson Educación.
\bibitem[SEP(2023)]{sep2023pensamiento} Secretaría de Educación Pública. (2023). \textit{Programa de estudios del recurso sociocognitivo Pensamiento Matemático I}. Subsecretaría de Educación Media Superior.
\bibitem[Star et~al.(2015)]{star2015algebra} Star, J. R., Caronongan, P., Foegen, A., Furgeson, J., Keating, B., Larson, M. R., Lyskawa, J., McCallum, W. G., Porath, J., y Zbiek, R. M. (2015). \textit{Teaching strategies for improving algebra knowledge}. IES.
\end{thebibliography}

\end{document}
